% =====================================================================
% INTRODUCTION
% =====================================================================
\section{مقدمه}
\label{sec:introduction}

\subsection{هدفِ یک ارزیابی امنیتی \en{LLM}}
\label{sec:intro:target}

یک مدل زبانی در محیط عملیاتی به‌ندرت مستقیماً در دسترس است. مدل پشت
برنامه‌ای می‌نشیند که تصمیم می‌گیرد چه چیزی وارد \en{context} آن شود،
اجازهٔ بازیابی چه چیزی را داشته باشد، با آنچه تولید می‌کند چه بشود و
خروجی‌اش مجاز به ایجاد چه اثرهای جانبی باشد: یک \en{input filter} پیش از
تولید، یک \en{retriever} روی پیکره‌ای که کاربران می‌توانند در آن بنویسند،
یک رمزگشا که تصویر بارگذاری‌شده را به متنِ \en{prompt} تبدیل می‌کند، یک
\en{output filter}، و یک لایهٔ \en{agent} که پرس‌وجوهای نوشته‌شده به‌دست
مدل را اجرا می‌کند. هر یک از این گام‌ها کدِ برنامه است، و هر یک جایی است
که دادهٔ تأمین‌شده به‌دست مهاجم می‌تواند اقتدارِ یک دستور را به دست آورد.

این همان مشاهدهٔ ساختاری پشت \en{indirect prompt injection}
است~\cite{greshake2023}: مدل نمی‌تواند به‌طور قابل اتکا دستوری را که به
او \emph{داده} شده از دستوری که به او \emph{نشان} داده شده تشخیص دهد،
چون تا وقتی هر دو به او می‌رسند همان توکن‌ها در همان \en{context}‌اند.
\en{Greshake} و همکاران نتیجه را چنین بیان می‌کنند که پردازشِ دادهٔ
بازیابی‌شدهٔ غیرقابل‌اعتماد قیاس‌پذیر با اجرای کد دلخواه
است~\cite[Sec.~2]{greshake2023}. از این‌جا برمی‌آید که پرسش امنیتیِ
شایسته دربارهٔ یک سامانهٔ مستقر، در وهلهٔ اول پرسشی دربارهٔ مدل نیست: مدلی
که به‌تنهایی کاویده شود نمی‌تواند خواندنِ میان‌مستأجریِ پایگاه‌داده را بروز
دهد، نمی‌تواند کوپنی را نشت دهد که یک بررسیِ \en{authorization} سمت سرور
هرگز آن را در \en{context} او قرار نمی‌داد، و نمی‌تواند از راه یک فرم نظر
مسموم شود. این شکست‌ها خواصِ یک خط لوله‌اند. با این حال بیشتر کارهای
منتشرشدهٔ \en{red teaming} مدل را واحد تحلیل می‌گیرند، و به دلیل خوبی هم:
مدل همان مؤلفه‌ای است که میان استقرارها تعمیم می‌یابد. پیامدش یک شکاف
است. سازمانی که یک برنامهٔ مبتنی بر \en{LLM} را اداره می‌کند انبوهی از
نتایج در سطح مدل در اختیار دارد و راهنمایی اندکی دربارهٔ اینکه چگونه
\emph{سامانهٔ خودش} را در برابر \emph{سیاست خودش} چنان بسنجد که مهندس
دیگری بتواند آن را بازپخش کند.

\subsection{تشخیص، بخش دشوار کار است}
\label{sec:intro:detection}

پرکردنِ آن شکاف مستقیماً به گامی می‌رسد که آزمون خودکار امنیت \en{LLM} نه
می‌تواند از آن بگذرد و نه می‌تواند تعمیمش دهد: تصمیم دربارهٔ اینکه آیا یک
خروجی شکست به شمار می‌آید. این تصمیم به استقرار وابسته است، چون یک رشتهٔ
یکسان در یک سامانه نقض سیاست است و در سامانهٔ دیگر پاسخی معمولی، و هیچ
\en{detector} قابل حملی نمی‌تواند به سیاست هیچ‌کدام حساس باشد. بنابراین
\en{detector}های عمومی به پرسشی عمومی پاسخ می‌دهند (\emph{آیا این شبیه
\en{jailbreak} است؟}) و سقفی از دقت را به ارث می‌برند که میزان شباهت یک
خروجی بی‌خطر به یک خروجی ناقض تعیینش می‌کند. ما این سقف را روی همین
استقرار اندازه می‌گیریم به‌جای آنکه فرضش کنیم (بخش~\ref{sec:res:agreement}).

ارزیابیِ یک برنامه می‌تواند بهتر عمل کند، مشروط به آنکه سیاستِ برنامه
آن‌قدر خوانا باشد که بتوان رونویسی‌اش کرد؛ و همین شرط است که مطالعهٔ حاضر
از آن بهره می‌برد. هدف، در کدِ خودش، دقیقاً همان واژه‌ای را نام می‌برد که
هر \en{persona} نباید بگوید؛ \en{PII} را با شکل‌های دقیقاً مشخص‌شده تولید
می‌کند؛ و به هر حسابِ کاشته‌شده یک \en{routing flag} حدس‌ناپذیر در
پایگاه‌دادهٔ ایزوله نسبت می‌دهد. هر یک از این ثابت‌ها به معنای دقیق کلمه
\en{ground truth} است: حضورش در یک خروجی نقض سیاست است و نبودش نه، بی‌آنکه
هیچ قضاوت شخصی در میان باشد. این کار یک مسئلهٔ امتیازدهیِ ذهنی را به
مسئله‌ای عینی تبدیل می‌کند، به بهای \en{detector}هایی که به هیچ سامانهٔ
دیگری منتقل نمی‌شوند؛ دقیقاً همان معامله‌ای که یک ارزیابی، در برابر یک
\en{benchmark}، باید بکند. از آن دو مسئله برمی‌خیزد. \en{ground truth}
باید به \en{detector} \emph{برسد}، چون مجموعه‌ای که پنج \en{persona} را در
یک فرایند جاروب می‌کند باید هر تلاش را در برابر رازِ \emph{همان تلاش}
امتیاز دهد. و یک \en{detector} مبتنی بر \en{ground truth} هنوز باید نقض
را \emph{بازشناسد}؛ رازی که یک \en{filter} ساده‌انگارانه آن را
\code{C-H-E-E-S-E} نوشته هنوز یک نشت است، و هر پالایشی که حساسیت می‌خرد،
از دقت هزینه می‌کند. هر دو در ادامه دغدغهٔ طراحیِ درجه‌یک‌اند و دومی به‌جای
فرض‌شدن اندازه‌گیری می‌شود.

\subsection{مشارکت‌ها}
\label{sec:intro:contributions}

این مقاله یک ارزیابی امنیتیِ سیاست‌آگاه و در سطح برنامه را از
\en{PwnzzAI Shop} گزارش می‌کند، یک برنامهٔ وبِ عمداً آسیب‌پذیرِ مبتنی بر
\en{LLM}، که تماماً به‌شکل \en{plugin}هایی برای \en{Garak}، پویشگر متن‌بازِ
امنیت \en{LLM}، ساخته شده
است.\footnote{\lr{\url{https://github.com/NVIDIA/garak}}. اینجا به‌عنوان
ابزاری نام برده می‌شود که ارزیابی در آن بیان شده، نه به‌عنوان منبعِ
ادعاها.} مجموعه، مصنوعات و مستنداتش سابقهٔ همراهِ این
مقاله‌اند~\cite{pwnzzrepo,pwnzzdocsoverview}.

\begin{enumerate}
\item \textbf{یک معماری \en{Garak}-بومی برای ارزیابی یک برنامه به‌جای یک
مدل} (بخش~\ref{sec:architecture}): نُه سطح که به‌شکل \en{generator}
بسته‌بندی و در زمان اجرا در فضای نام و \en{cache} خودِ \en{Garak} ثبت
می‌شوند بی‌آنکه فایلی در محل نصب کپی شود، به‌طوری که هر اندازه‌گیری از راه
\en{CLI} استاندارد بازپخش می‌شود و هستهٔ \en{Garak} دست‌نخورده می‌ماند.

\item \textbf{تشخیصِ مبتنی بر \en{ground truth} روی یک کانال جانبیِ پاسخ}
(بخش‌های~\ref{sec:arch:notes} تا~\ref{sec:arch:detectors}): دوازده
\en{detector} که یک نقض مشخص را در برابر ثابت‌های رونویسی‌شده از کدِ
پین‌شده و بازوارسی‌شده با آزمون‌های قراردادی می‌آزمایند، با
\en{ground truth} هر تلاش که در \code{Message.\bk notes} سفر می‌کند، و
امتناع (\code{None}، هرگز $0.0$) هرگاه \en{detector} نتواند قضاوت کند.

\item \textbf{یک \en{LLM-as-a-judge} کالیبره‌شده به‌عنوان نظر دومِ
اختیاری} (بخش~\ref{sec:arch:judge})، برای پوششِ آنچه \en{ground truth}
نمی‌بیند: پاسخی که مهاجم را راضی می‌کند بی‌آنکه با الگو تطبیق یابد. این
داور به‌صورت پیش‌فرض خاموش است، هرگز اصلی نیست، و کالیبراسیونِ ترتیب
\en{schema} آن به‌جای فرض‌شدن گزارش می‌شود.

\item \textbf{طراحیِ لایه‌جداکننده با کنترل‌های درون‌ساخت}
(بخش~\ref{sec:methodology}): یک پیکرهٔ حملهٔ ثابت که روی پنج سطح
\en{persona} و ده پیکربندی \en{guardrail} جاروب می‌شود که تنها در لایهٔ
دفاعیِ حاضر تفاوت دارند؛ \en{poisoning} که در برابر یک کنترلِ سالمِ
جفت‌شده و برازش‌شده در همان اجرا سنجیده می‌شود؛ یک \en{mitigation} در
\en{retrieval} که با پیکره و \en{prompt}های ثابت روشن و خاموش می‌شود؛ و
یک بررسیِ سلامتِ تحویل که تلاش‌های غیرمستقیمِ بدونِ رسیدنِ \en{payload} را
کنار می‌گذارد.

\item \textbf{نتایج روی \num{678} تلاش ارزیابی‌شده}
(بخش~\ref{sec:results}): تمرکز موفقیت در تکنیک‌هایی که هرگز هدف خود را
نام نمی‌برند؛ دو \en{guardrail} از ده که بهتر از هیچ نیستند؛ آستانه‌ای در
\en{poisoning} که چهار گامِ بودجه از آغاز اثر عقب می‌ماند؛ و یک
\en{detector} عمومی با \en{recall} کامل و \en{precision} برابر
\num{0.293}.

\item \textbf{یک مطالعهٔ \en{ablation}} (بخش~\ref{sec:ablation}) که سهم
هر لایهٔ دفاعی را از سهم هر جزء دستگاه اندازه‌گیری جدا می‌کند. از ده
مؤلفه، سه مورد تضمین‌هایی‌اند که این اجراها هرگز به‌کارشان نبردند، پنج
مورد یک عدد گزارش‌شده را تغییر می‌دهند (یکی با ضریب \num{3.4})، و دو مورد
به‌جای تضعیف یک اندازه‌گیری، آن را از میان برمی‌دارند.

\item \textbf{راهکارهای متصل به شواهد} (بخش~\ref{sec:disc:mitigations})،
که هر یک به سیگنالِ \en{detector} پشتیبانش و به لایه‌ای از خط لوله که روی
آن عمل می‌کند گره خورده است، همراه با دو ردهٔ مثبت کاذب که علیه
\en{detector}های خودمان گزارش شده‌اند.
\end{enumerate}
